\chapter{名词变格：\\\Large\itshape\textcolor{gray}{三性四格二数}}
\label{chap:noun}
\section{性}
德语与英语最主要区别之一在于，不仅人有性，物也有性。\textcolor{gray}{（英语物也有性的例外是：对船、车的称呼用she）}
要明白性是德语名词不可分割的一部分。要注意的一点是，语法性与天然性别可能不同，如(das Mädchen)。与之对应的，人称代词也可以用来指代物。 \label{sec:noun:gender} \\
\begin{myquote}
    背名词，必带冠。
\end{myquote}

\section{格}
\subsection{格的顺序}
\begin{itemize}
    \item 数字顺序：主属与宾，NGDA\textcolor{gray}{（语法书常用到第几格的说法）}
    \item 课本顺序：主宾与属，NADG
\end{itemize}
\subsection{变格规律}
\label{subsec:noun:declension}
名词的变格相对简单，黄书第一章就是它。
单数名词只有阳中的属格加-(e)s（多音节-s,单音节-es）。复数名词与格加-n。阳性名词中有部分“弱名词”，除单数主格外一律加-(e)n。
\textcolor{gray}{（复与n,这个在冠词、人称代词还会见到！所以说德语有规可循！）}
\begin{myquote}
    复与-n，阳弱-(e)n，单阳中属-(e)s。
\end{myquote}

\subsubsection{阳性弱变化名词}
阳弱三类
\begin{itemize}
    \item -e 结尾
          \begin{center}
              \begin{tabular}{|c|c|}
                  \hline
                  \rowcolor{gray}
                  word       & meaning      \\
                  \hline
                  Junge      & boy          \\
                  Kollege    & colleague    \\
                  Kunde      & customer     \\
                  Neffe      & nephew       \\
                  Zeuge      & witness      \\
                  \hline
                  Chinese    & chinese      \\
                  Franzose   & french       \\
                  Russe      & russian      \\
                  \hline
                  Affe       & ape          \\
                  Hase       & hare         \\
                  Löwe       & lion         \\
                  \hline
                  Biologe    & biologist    \\
                  Psychologe & psychologist \\
                  \hline
              \end{tabular}
          \end{center}
    \item -ant, -ent, -ist, -at 结尾
          \begin{center}
              \begin{tabular}{|c|c|}
                  \hline
                  \rowcolor{gray}
                  word                   & meaning              \\
                  \hline
                  Elefant                & elephant             \\
                  \hline
                  Assistent              & assistant            \\
                  Student                & (university) student \\
                  Präsident              & president            \\
                  Dozent                 & lecturer             \\
                  \hline
                  \underline{Jour}nalist & journalist           \\
                  Polizist               & police               \\
                  Optimist               & optimist             \\
                  Pessimist              & pessimist            \\
                  \hline
                  Diplomat               & diplomat             \\
                  \hline
              \end{tabular}
          \end{center}
    \item 部分人或生物
          \begin{center}
              \begin{tabular}{|c|c|}
                  \hline
                  \rowcolor{gray}
                  word    & meaning   \\
                  \hline
                  Bär     & bear      \\
                  \hline
                  Bauer   & farmer    \\
                  Herr    & Mister    \\
                  Mensch  & person    \\
                  Nachbar & neighbour \\
                  \hline
              \end{tabular}
          \end{center}
\end{itemize}


\tikzset{
    nountable/.style={
            matrix of nodes,
            nodes={
                    anchor=center,
                    minimum width=3em,
                    minimum height=2.5em,
                    draw=black,
                },
            draw=black,
            inner sep=0,
            outer sep=0,
            node distance=0,
        },
}
\subsection{格的定义、疑问代词}
\label{subsec:noun:inqpron}
依赖句中角色，由动/介决定，
\begin{center}
    \begin{tabular}{|c|c|c|}
        \hline
        \rowcolor{gray}
        格 & 定义 & 疑       \\
        \hline
        N & 主语 & wer/was \\
        \hline
        A & 直宾 & wen/was \\
        \hline
        D & 间宾 & wem     \\
        \hline
        G & 所属 & wessen  \\
        \hline
    \end{tabular}
\end{center}
格区分名词性成分所属句子成分。因此德语句子顺序相对灵活。e.g.:
\begin{itemize}
    \item Der Mann liebt die Frau.
    \item Die Frau liebt der Mann.
\end{itemize}
\noindent
大部分带宾语动词是第四格。 \\
关于双宾：顺序是，名与前，代宾前。区别是，人间物直，静三动四。 \\
关于介宾搭配：疑问词用wo+prep，或者woher/wohin(介词不确定时)
\subsubsection{wo}
\begin{center}
    \begin{tabular}{|c|c|c|c|}
        \hline
        \rowcolor{gray}
        wh                     & prep & obj case                                 & example                                        \\
        \hline
        \multirow{2}{*}{wo}    & in   & \pill[border=true, shape={type=olive}]   & Wo bist du? Ich bin in der Schweiz.            \\
        \cline{2-4}
                               & auf  & \pill[border=true, shape={type=olive}]   & Wo kauft er ein? Er Kauft auf dem Markt.       \\
        \hline
        woher                  & aus  & \pill[border=true, shape={type=olive}]   & Woher kommst du? Ich komme aus der Schweiz.    \\
        \hline
        \multirow{3}{*}{wohin} & nach & \pill[border=true, shape={type=olive}]   & Wohin fliegst du? Ich fliege nach der Schweiz. \\
        \cline{2-4}
                               & in   & \pill[border=true, shape={type=diamond}] & Wohin fliegst du? Ich fliege in die Schweiz.   \\
        \cline{2-4}
                               & auf  & \pill[border=true, shape={type=diamond}] & Wohin geht er? Er geht auf den Markt.          \\
        \hline
    \end{tabular}
\end{center}
\subsection{名词数，疑问代词}
\begin{center}
    \begin{tabular}{|c|c|c|}
        \hline
        \rowcolor{gray}
        数词         & 疑问词       \\
        \hline
        uncoutable & wie viel  \\
        \hline
        countable  & wie viele \\
        \hline
    \end{tabular}
\end{center}
z.B.:\\
Wie viele Übungen machen Sie heute? \\
注意Wie Viele/Viel \sqr[border=true] 算作一个句子成分。

\subsection{可视化格、性}
用颜色形状表示性、格，数算一种特殊性。
\begin{center}
    \begin{tikzpicture}
        \matrix (m1) [nountable,
            row 1/.style={nodes={minimum height=1.5em, fill=gray!50}},
        ]{
            格 & 形                    & 性 & 色                     \\
            N & \sqr[border=true] & M & \sqr[color=green]  \\
            A & \diam[border=true]   & N & \sqr[color=purple] \\
            D & \olive[border=true]  & F & \sqr[color=yellow] \\
            G & \tri[border=true]    & P & \sqr[color=red]    \\
        };
        \node(l1) [above=0.03cm of m1] {\small{Table 1-1 可视化}};
    \end{tikzpicture}
\end{center}


\newgeometry{left=1cm, right=1cm, top=1cm, bottom=1cm}
\section{名词性、数规律}
\label{subsec:noun:gender:number}
使用方法：由上到下看图。
\begin{itemize}
    \item 先看词形，若匹配后缀，可推断对应的性、数（有5个能直接确定）
    \item 若无匹配，再看所属词族，若匹配，则可猜出性
    \item 根据性和音节数，便可以猜出复数形式
    \item 此外合成词的性取决于尾词
    \item 变音有四：a\textrightarrow{ä},o\textrightarrow{ö},u\textrightarrow{ü},au\textrightarrow{äu}
    \item 小结论：-(e)n和-s从无变音，单阴-e总是变音，-er从无阴性，-
\end{itemize}
\begin{myquote}
    形义推性，性音推数，否则特殊，合成看尾。
\end{myquote}
\tikzset{
    glossary/.style={
            nodes={
                    minimum height=0.5cm,
                    minimum width=4cm,
                    text width=4cm,
                    font=\linespread{0.9}\bfseries,
                    inner sep=0,
                    outer sep=0,
                    thick
                },
            node distance=0.1cm,
            align=center,
        },
    label/.style={
            anchor=center,
        },
    masculine/.style={
            draw=green,
            text=green,
        },
    neuter/.style={
            draw=purple,
            text=purple,
        },
    feminine/.style={
            draw=yellow,
            text=yellow,
        },
    plural/.style={
            text=red,
            align=left,
            minimum width=1cm,
            text width=1cm,
        },
    note/.style={
            text=gray,
            align=left,
            minimum width=1.5cm,
            text width=1.5cm,
        },
    ambiguous/.style={
            draw=darkgray,
            text=darkgray,
        },
    fitbox/.style={
            inner sep=0.1cm,
            outer sep=0.1cm,
            rounded corners=0.1cm,
        },
    proportionarrow/.style={
            -latex,
            shorten >=0,
            shorten <=0
        },
    arrowtext/.style={
            font=\sffamily,
            text=gray,
        },
}
\begin{tikzpicture}[glossary]
    % prefixes with >99% certainty
    \node(m1) [masculine] {-ant, -and};
    \node(m2) [masculine, below=of m1] {-ast, -ling, -ich, -g};
    \node(m3) [masculine, below=of m2] {-or, -us, -ist, -ismus};
    \node(p1) [plural, right=0.3cm of m1] {$\text{-en}^*$};
    \node(p2) [plural, right=0.3cm of m2] {-e};
    \node(t1) [masculine, fitbox, fit=(m1) (m2) (m3) (p1) (p2)] {};
    \node(l1) [label, text=green, above=1cm of t1] {\Large{M}};

    \node(n1) [neuter, right=0.6cm of p1] {-chen, -lein};
    \node(n2) [neuter, below=of n1] {-il};
    \node(n3) [neuter, below=of n2] {-icht, -it, -ma, -ment,\\-tum, -tel, -um};
    \node(n4) [neuter, below=of n3] {Ge-};
    \node(p3) [plural, right=0.3cm of n1] {$\text{-}^*$};
    \node(p4) [plural, right=0.3cm of n2] {-e};
    \node(t2) [neuter, fitbox, fit=(n1) (n2) (n3) (n4) (p3) (p4)] {};
    \node(l2) [label, text=purple] at ($(t2.center)!(l1.center)!(t2.center)$) {\Large{N}};

    \node(f1) [feminine, right=0.6cm of p3] {-ei, -heit, -keit,\\-ung, -schaft,\\-in, -frau};
    \node(f2) [feminine, below=of f1] {-a -ik, -ur,\\-sion, -tion,\\-anz, -enz,\\-tät, -sis, -ie};
    \node(f3) [note, right=0.3cm of f2] {Foreign};
    \node(t3) [feminine, fitbox, fit=(f1) (f2) (f3)] {};
    \node(l3) [label, text=yellow] at ($(t3.center)!(l1.center)!(t3.center)$) {\Large{F}};

    \node(tt) [draw=none, fit=(t1) (t2) (t3)] {};

    % % ambiguous prefixes
    \node(r1) [draw=none] at ([yshift=-0.9cm]$(l1.center)!(tt.south)!(l1.center)$) {};
    \node(a1) [ambiguous] at ([xshift=-0.7cm]r1.south) {-en, -el, -r};
    \draw[proportionarrow, masculine] (a1) -- (t1) node[arrowtext, midway] {\small{60\%}};

    \node(r2) [draw=none] at ([yshift=-0.9cm]$(l2.center)!(tt.south)!(l2.center)$){};
    \node(a2) [ambiguous] at ([xshift=2cm]r2.south) {-nis, -sal};
    \node(p5) [plural, right=0.3cm of a2] {-e};
    \draw[proportionarrow, neuter] (a2) -- (t2) node[arrowtext, midway] {\small{70\%}};
    \draw[proportionarrow, feminine] (a2) -- (t3) node[arrowtext, midway] {\small{30\%}};

    \node(r3) [draw=none] at ([yshift=-2.4cm]$(l2.center)!(tt.south)!(l2.center)$) {};
    \node(a3) [ambiguous] at ([xshift=-0.5cm]r3.south) {-al, -an, -ar, -är, -at,\\-ent, -ett,\\-ier, -iv, -o, -on};
    \draw[proportionarrow, masculine] (a3) -- (t1) node[arrowtext, midway] {\small{10\%}};
    \draw[proportionarrow, neuter] (a3) -- (t2) node[arrowtext, midway] {\small{90\%}};

    \node(r4) [draw=none] at ([yshift=-0.9cm]$(l3.center)!(tt.south)!(l3.center)$) {};
    \node(a4) [ambiguous] at ([xshift=1cm]r4.south) {-e};
    \draw[proportionarrow, feminine] (a4) -- (t3) node[arrowtext, midway] {\small{90\%}};

    \node(t4) [draw=darkgray, densely dashed, fitbox, fit=(a1) (a2) (a3) (a4)] {};

    \node(bt1) [draw=darkgray, fitbox, fit=(t1) (t2) (t3) (t4)]{}; % big fit box 1

    % meaning groups
    \node(mg1) [masculine] at ([yshift=-1cm]$(r1.center)!(t4.south)!(r1.center)$) {Cars};
    \node(mg2) [masculine, below=of mg1] {Currency};
    \node(mg3) [masculine, below=of mg2] {Days, Months,\\Seasons};
    \node(mg4) [masculine, below=of mg3] {Drinks};
    \node(mg5) [masculine, below=of mg4] {Weather};
    \node(mg6) [masculine, below=of mg5] {Mountains};
    \node(mg7) [masculine, below=of mg6] {Space};
    \node(mg8) [masculine, below=of mg7] {Minerals};
    \node(t5) [masculine, fitbox, fit=(mg1) (mg2) (mg3) (mg4) (mg5) (mg6) (mg7) (mg8)] {};

    \node(ng1) [neuter] at ($(r2.center)!(mg1.center)!(r2.center)$) {Cars};
    \node(ng2) [neuter, below=of ng1] {Countries};
    \node(ng3) [neuter, below=of ng2] {Continents, Cities,\\Provinces};
    \node(ng4) [neuter, below=of ng3] {Hotel, Cafe};
    \node(ng5) [neuter, below=of ng4] {Restaurant, Theatres};
    \node(ng6) [neuter, below=of ng5] {Languages};
    \node(ng7) [neuter, below=of ng6] {Notations};
    \node(ng8) [neuter, below=of ng7] {Colors};
    \node(ng9) [neuter, below=of ng8] {Elements};
    \node(ng10) [neuter, below=of ng9] {Units};
    \node(t6) [neuter, fitbox, fit=(ng1) (ng2) (ng3) (ng4) (ng5) (ng6) (ng7) (ng8) (ng9) (ng10)] {};

    \node(fg1) [feminine] at ($(r4.center)!(mg1.center)!(r4.center)$) {Numerals};
    \node(fg2) [feminine, below=of fg1] {Tree, Fruits,\\Flowers};
    \node(fg3) [feminine, below=of fg2] {Airplane, Motorcycle,\\Ships};
    \node(t7) [feminine, fitbox, fit=(fg1) (fg2) (fg3)] {};

    \node(a5) [ambiguous] at ([xshift=1cm, yshift=-2cm]t5.south) {Persons, Rivers, ...};
    \node(a6) [ambiguous, right=1cm of a5] {Animals, ...};
    \node(a7) [ambiguous, right=1cm of a6] {...};
    \node(t7) [draw=darkgray, densely dashed, fitbox, fit={($(a5.south)!(t5.west)!(a6.south)$) ($(a5.north)!(t7.east)!(a6.north)$)}] {};

    \node(bt2) [draw=darkgray, fitbox, fit={($(t4.west)!(t7.south)!(t4.west)$) ($(t4.east)!(t5.north)!(t4.east)$)}]{}; % big fit box 2

    % syllables
    \node(pp1) [draw=red, text=red] at ([yshift=-1.3cm]$(r1.center)!(t7.south)!(r1.center)$) {\textcolor{gray}{monosyllabic:} {\freesans{\char"2E1A}}e};
    \node(pp2) [draw=red, text=red] at ($(r2.center)!(pp1.center)!(r2.center)$) {{\textcolor{gray}{monosyllabic:} \freesans{\char"2E1A}}er};
    \node(tt1) [draw=none, below=0.3cm of pp1] {};
    \node(tt2) [draw=none, below=0.3cm of pp2] {};
    \node(pp3) [draw=red, text=red, fit=(tt1) (tt2)] {\raisebox{-1ex}{\textcolor{gray}{polysylabic:}-e}};

    \node(tt3) [draw=none, below=0.3cm of tt1] {};
    \node(tt4) [draw=none, below=0.3cm of tt2] {};
    \node(pp4) [draw=red, text=red, fit=(tt3) (tt4)] {\raisebox{-1ex}{\textcolor{gray}{-a, -e, -i, -o, -u, -y}-s}};

    \node(tt5) [draw=none] at ($(r4.center)!(pp1.center)!(r4.center)$) {};
    \node(tt6) [draw=none, below=0.3cm of tt5] {};
    \node(pp5) [draw=red, text=red, fit=(tt5) (tt6)] {\raisebox{-1ex}{-n\textcolor{gray}{ (weak nouns)}}};

    \node(bt3) [draw=darkgray, minimum height=2.5cm, fitbox, fit={($(t4.west)!(pp4.south)!(t4.west)$) ($(t4.east)!(pp1.north)!(t4.east)$)}]{}; % big fit box 3
\end{tikzpicture}
\restoregeometry
\newpage
特殊情况包括弱名词，特殊复数。